\section{The static input and its normalization}
\label{sec:static-input}

We record the precise input from the companion article
\cite[Theorem~2.2]{StaticTrees}. Its constants and small exponents
belong to one fixed limiting tree. No uniform exponent across different
trees is needed. The elementary functional identities are included here
to fix the normalizations used in the flow calculation.
The remaining static proof, including path selection and the augmented
elliptic construction, is supplied by the complete companion.

\subsection{The geometric data and the functional}
\label{ell:setup}

Let $\mathcal T$ be a finite rooted tree of four-dimensional Ricci-flat
spaces $(V_v,h_v)$.  A vertex has one exterior end and one punctured end
for each child.  Every end has Cartesian coordinates on a Euclidean
annulus modulo a finite subgroup of $O(4)$ acting freely on $S^3$.
The group on a joined end is nontrivial.  The two groups on an edge
are identified; a fixed orthogonal identification is included in the
charts.  Put $\tau=\log r$ on an exterior end and $\tau=-\log r$ at a
puncture.  In these charts, in vertex units,
\begin{equation}
 h_v=\delta+O(e^{-\mu\tau})
 \label{ell:end-decay}
\end{equation}
on joined ends, with derivatives at scale, for some $\mu>0$.
On the unjoined exterior end of the root we require instead
$h_v-\delta=O(r^{-q})$, where $q>1$.
All remaining parts are smooth compact manifolds with the indicated
boundary collars.  Bounds through order $20$ suffice in the data below.

For an edge $e$ from $p$ to $c$, truncate both end coordinates at
$\tau=L_e$ and identify the boundary collars by
\begin{equation}
 x_p=e^{-2L_e}x_c,\qquad a_c=a_p e^{-2L_e},\qquad a_{\rm root}=1.
 \label{ell:scales}
\end{equation}
Thus the metric tensor in vertex units is multiplied by $a_v^2$ in
physical units.  The physical radius $\rho=a_vr$ agrees in both charts,
and $\tau_p+\tau_c=2L_e$.  Gluing just the boundary faces, with their
compatible smooth collars, gives the same manifold.

Consider a sequence of metrics $g$ for which every $L_e\to\infty$,
which converges to $h_v$ on each compact subset of every open vertex.
Assume the same end estimates as in Equation~\eqref{ell:end-decay}
on each joining half, and the stated root estimate, with uniform
finite derivative bounds.  On a compact core put $\rho=a_v$.
Suppose $M\to0$ and
\begin{equation}
 |\rho^2\Ric(g)|_j\le CM
 \begin{cases}
 e^{-\mu\tau},&\text{on joined half-ends},\\
 r^{-q},&\text{on the root exterior end},\\
 1,&\text{on vertex cores},
 \end{cases}
 \qquad j=10.
 \label{eq:tree-ricci-error}
\end{equation}
Here the tensor magnitude is physical, and the subscript includes
covariant derivatives multiplied by the corresponding powers of
$\rho$.  The metric bounds convert these to coordinate estimates.
Decreasing $\mu$ or $q$, while retaining $\mu>0$ and $q>1$, is allowed.

Define, in physical root units,
\begin{equation}
 E(g)=\lim_{D\to\infty}\left\{
 \int_{r<D}R(g)\,d\mu_g-
 \int_{r=D}(\partial_jg_{ij}-\partial_ig_{jj})n_i\,dS_\delta
 \right\}.
 \label{eq:energy-functional}
\end{equation}
All chart integrals are quotient integrals.  The two terms inside the
braces need not converge separately.

The boundary subtraction is the mass term in the renormalized
Einstein--Hilbert expression discussed by Deruelle and Ozuch
\cite[Definition~4.2, Remark~4.3, and Proposition~4.4]{DeruelleOzuch2025}.
Their ALE version of Perelman's functional also minimizes over an
auxiliary potential; the functional $E$ here is the unminimized
scalar-curvature-minus-flux expression.

\begin{lemma}[Renormalization and first variation]
\label{ell:functional}
The functional in Equation~\eqref{eq:energy-functional} is defined on
the preceding metrics.  For a variation with the same root decay,
\begin{equation}
 DE_g(\dot g)=\int\left\langle\tfrac12R(g)g-\Ric(g),\dot g\right\rangle_g
 \,d\mu_g.
 \label{ell:first-variation}
\end{equation}
It is invariant under diffeomorphisms isotopic to the identity whose
root displacement, including scaled derivatives throughout the
isotopy, is $O(r^{1-q'})$ for some $q'>1$.  A global length scaling
by $s$, accompanied by a coordinate dilation at infinity restoring
the leading tensor $\delta$, multiplies $E$ by $s^2$.
Under Equation~\eqref{eq:tree-ricci-error}, $|E(g)|\le CM$.
\end{lemma}

\begin{proof}
Write $g=\delta+p$ on the root end.  The scalar-curvature density is
\[
 R(g)\sqrt{\det g}=
 \partial_i(\partial_jp_{ij}-\partial_ip_{jj})+
 O(|p||\partial^2p|+|\partial p|^2).
\]
The last expression is $O(r^{-2q-2})$ and its radial integral is
bounded by $C\int_{r_0}^\infty r^{1-2q}\,dr$.  This proves convergence.
The linear boundary term in the variation of scalar curvature
cancels the variation of the subtracted integral.  Every remaining
boundary product is $O(D^{2-2q})$; it tends to zero.  Integration by
parts therefore gives Equation~\eqref{ell:first-variation}.
For a Lie derivative, the contracted Bianchi identity makes the
interior divergence vanish.  The remaining boundary term vanishes
with the stated displacement decay.  Integration along the isotopy
gives invariance.  Changing variables under a dilation in
Equation~\eqref{eq:energy-functional} gives the scaling law.

Differentiate this law at $s=1$.  Choose the compensating dilation
to be the identity inside a fixed outer root region.  Its metric
variation equals $2g$ on the interior and decays like $r^{-q}$ on
the root tail.  Its pairing with Ricci on the root is integrable
by $q>1$.  On a punctured half the interior estimate costs
\[
 CM a_p^2\int_{\tau_0}^{L_e}e^{-(2+\mu)\tau}\,d\tau\le CM a_p^2.
\]
On a child exterior half the corresponding quantity is
\[
 CM a_c^2\int_{\tau_0}^{L_e}e^{(2-\mu)\tau}\,d\tau
 \le CM a_p^2(1+L_e)e^{-\min(2+\mu,4)L_e},
\]
with a harmless change of constant for the fixed lower endpoint.
The finite sum of core contributions is bounded as well, since
$a_v\le1$.  Equation~\eqref{ell:first-variation} now gives
$2|E(g)|\le CM$.  In particular $E=0$ when $\Ric=0$.
\end{proof}

\begin{theorem}[Static tree inequality {\cite[Theorem~2.2]{StaticTrees}}]
\label{thm:tree-inequality}
For the fixed tree and the sequences just specified,
\begin{equation}
                         E(g)=o(M).
 \label{ell:little-oh}
\end{equation}
For $M=0$ the assertion means $E(g)=0$.
\end{theorem}

The companion proof constructs a finite-dimensional analytic family with
complex length parameters.  It does not assume that Einstein
deformations are unobstructed, or that the given metrics or the
coefficients of a fixed core metric are spatially analytic.

The theorem is applied in Proposition~\ref{final:functional-smallness}.
There we verify every hypothesis after extracting a fixed tree from an
arbitrary sequence of active times and material collars. The root
exterior has decay exponent greater than one because it is replaced
by the extension of Section~\ref{sec:collar}; the true joining necks
retain their own positive decay exponent.
